% Full proof completion for FLRV1-U00869; original theorem routing is lead-owned.
\section{Kripke soundness, prime theories, and strong completeness}
\label{sec:p08-ipl-kripke-completeness}

The common-world consequence relation of
Section~\ref{sec:intuitionistic-logic-01} has a complete proof by prime
theories. Its finite countermodels require a further construction, given
below. The semantic starting point is Kripke's persistent interpretation
\cite[pp.~93--95]{Kripke1965}. The proof here works directly with the
propositional rules of \hyperref[sheet:lj]{\(\mathsf{LJ}\)}.

\begin{hbconvention}{def:p08-ipl-completeness-setting}{Language, premise sets, and metatheory}
Fix a set \(\mathsf{Prop}\) of propositional atoms and let \(\mathcal F\)
be the set of finite formulas generated by
\[
  A ::= p\mid\bottom\mid(A\land A)\mid(A\lor A)\mid(A\to A),
  \qquad p\in\mathsf{Prop}.
\]
Negation abbreviates \(\neg A:=A\to\bottom\). The letters \(A,B,C,D\)
denote formulas; \(\Gamma,\Delta\) denote arbitrary subsets of
\(\mathcal F\), and subscripted contexts such as \(\Gamma_0\) denote
finite multisets when they occur in a sequent. Write
\(\Gamma\derives_{\mathsf{IPL}}A\) when some finite multiset
\(\Gamma_0\) of members of \(\Gamma\) has a derivation of
\(\Gamma_0\Rightarrow A\) in the propositional part of
\(\mathsf{LJ}\). This part contains Id, \(\bottom L\), WL, CL,
\(\land L,\land R,\lor L,\lor R_1,\lor R_2,\to L,\to R\),
Cut, and Mix. Its succedent contains exactly one formula; exchange is built
into the multiset convention. There are no quantifier or equality rules in
this propositional presentation.

The general completeness theorem is proved in classical ZFC. In particular,
Zorn's lemma supplies prime extensions for an arbitrary set
\(\mathsf{Prop}\); no countability hypothesis on this set or on
\(\Gamma\) is imposed. If an enumeration of \(\mathcal F\) is supplied,
the separate construction in Lemma~\ref{thm:p08-ipl-enumerated-extension}
works in classical ZF without invoking a choice principle. Neither
construction assumes an effective membership test for \(\Gamma\).
\end{hbconvention}

\begin{hbdefinition}{def:p08-ipl-completeness-semantics}{Local consequence and rooted frames}
A model is \(\mathcal M=(W,\leq,V)\), with a nonempty partially ordered
world set \(W\), a persistent valuation
\(V:\mathsf{Prop}\to\mathcal P(W)\), and the forcing clauses of
Section~\ref{sec:intuitionistic-logic-01}. The letters \(u,v,w\) denote
worlds of this model. Thus \(w\forces p\) means
\(w\in V(p)\), no world forces \(\bottom\), conjunction and
disjunction are evaluated at the same world, and
\[
 \mathcal M,w\forces A\to B
 \quad\Longleftrightarrow\quad
 \text{for every }v\geq w,
 \ \mathcal M,v\forces A\text{ implies }\mathcal M,v\forces B.
\]
Write \(\mathcal M,w\forces\Gamma\) if \(w\) forces every member of
\(\Gamma\). The relation \(\Gamma\entails_{\mathsf{IPL}}A\) means
that every such model and world forcing \(\Gamma\) also forces \(A\).
All premises are evaluated together at that world. A frame is
\emph{rooted at \(r\)} when \(r\in W\) and \(r\leq w\) for every
\(w\in W\).

There is also a preorder presentation, with reflexivity and transitivity
but no antisymmetry requirement, as in
Definition~\ref{def:p08-ipl-kripke-structure}. The soundness proof applies
to that presentation too. The canonical and finite countermodels below
are partial orders, and hence establish completeness for each of the two
frame classes. Allowing empty frames adds no counterexample to local
consequence, because an empty frame has no evaluation world.
\end{hbdefinition}

\begin{hblemma}{thm:p08-ipl-completeness-persistence}{proved=annex:p8.ipl.complete.persistence}{Persistence for the full propositional language}
In every model on a preorder, for every \(A\in\mathcal F\),
\[
 w\leq v\text{ and }\mathcal M,w\forces A
 \quad\Longrightarrow\quad \mathcal M,v\forces A.
\]
Consequently, forcing an arbitrary premise set persists to future worlds.
\end{hblemma}

\begin{hbtheorem}{thm:p08-ipl-completeness-soundness}{proved=annex:p8.ipl.complete.soundness}{Soundness of every propositional LJ rule}
Every derivable propositional sequent \(\Gamma_0\Rightarrow A\) is
valid at every world of every persistent preorder model. Therefore, for
every premise set \(\Gamma\subseteq\mathcal F\),
\[
 \Gamma\derives_{\mathsf{IPL}}A
 \quad\Longrightarrow\quad \Gamma\entails_{\mathsf{IPL}}A.
\]
The same assertion holds when consequence quantifies over preorder models.
The proof includes the displayed \(n\)-fold Mix rule for every
\(n\geq1\).
\end{hbtheorem}

\subsection{Finite support and prime extension}
\label{sec:p08-ipl-prime-extension}

\begin{hbdefinition}{def:p08-ipl-prime-theory}{Deductive closure and prime theories}
For \(\Gamma\subseteq\mathcal F\), define
\[
 \operatorname{Cn}_{\mathsf{IPL}}(\Gamma)
   :=\{A\in\mathcal F:\Gamma\derives_{\mathsf{IPL}}A\}.
\]
A \emph{theory} \(T\) is a subset of \(\mathcal F\) closed under this
derivability relation. It is \emph{consistent} when
\(T\not\derives_{\mathsf{IPL}}\bottom\), and it is a
\emph{prime theory} when it is consistent and, for every \(A,B\),
\[
 A\lor B\in T\quad\Longrightarrow\quad A\in T\text{ or }B\in T.
\]
For a theory, membership in \(T\) is equivalent to derivability from
\(T\). In particular, the disjunction condition concerns disjunctions
derivable from \(T\). A set \(S\) \emph{avoids \(C\)} when
\(S\not\derives_{\mathsf{IPL}}C\). It is \emph{maximal avoiding
\(C\)} when it avoids \(C\) and every proper extension inside
\(\mathcal F\) derives \(C\).
\index{prime theory!intuitionistic}
\end{hbdefinition}

\begin{hblemma}{thm:p08-ipl-finite-support-operations}{proved=annex:p8.ipl.complete.finite-support}{Finite support, substitution of derivations, and deduction}
Finite-support derivability has the following properties for arbitrary
premise sets.
\begin{enumerate}
\item If \(A\in\Gamma\), then \(\Gamma\derives_{\mathsf{IPL}}A\).
  Enlarging \(\Gamma\) preserves derivability, and any derivation uses
  only finitely many members of \(\Gamma\).
\item If \(\Delta\derives_{\mathsf{IPL}}C\) and
  \(\Gamma\derives_{\mathsf{IPL}}D\) for every \(D\in\Delta\),
  then \(\Gamma\derives_{\mathsf{IPL}}C\).
  Thus \(\operatorname{Cn}_{\mathsf{IPL}}\) is extensive, monotone,
  and idempotent.
\item The deduction equivalence is
  \[
   \Gamma\cup\{A\}\derives_{\mathsf{IPL}}B
   \quad\Longleftrightarrow\quad
   \Gamma\derives_{\mathsf{IPL}}A\to B.
  \]
\item Derivability admits conjunction introduction and each projection,
  each disjunction introduction, modus ponens, and explosion. It also
  admits reasoning by cases: if \(\Gamma\) derives \(A\lor B\),
  \(\Gamma\cup\{A\}\) derives \(C\), and
  \(\Gamma\cup\{B\}\) derives \(C\), then \(\Gamma\) derives
  \(C\).
\end{enumerate}
All these assertions refer to the exact single-succedent calculus and
finite-multiset contexts on \hyperref[sheet:lj]{the LJ sheet}.
\end{hblemma}

\begin{hblemma}{thm:p08-ipl-maximal-is-prime}{proved=annex:p8.ipl.complete.maximal-prime}{Maximal avoidance gives a prime theory}
Every set maximal avoiding a formula \(C\) is a prime theory and does
not contain \(C\). More generally, if \(S\) avoids \(C\) and
\(S\derives_{\mathsf{IPL}}A\lor B\), then at least one of
\(S\cup\{A\}\), \(S\cup\{B\}\) avoids \(C\).
\end{hblemma}

\begin{hblemma}{thm:p08-ipl-prime-extension}{proved=annex:p8.ipl.complete.prime-extension}{Prime extension for arbitrary premise sets}
In classical ZFC, if \(\Gamma\not\derives_{\mathsf{IPL}}C\),
there is a prime theory \(T\supseteq\Gamma\) with \(C\notin T\).
The language and premise set may have arbitrary set cardinality.
Taking \(C=\bottom\) shows that every consistent premise set extends
to a prime theory.
\end{hblemma}

\begin{hblemma}{thm:p08-ipl-enumerated-extension}{proved=annex:p8.ipl.complete.enumerated-extension}{Prime extension from a supplied formula enumeration}
Suppose \((D_n)_{n\in\mathbb N}\) enumerates every member of
\(\mathcal F\), and \(\Gamma\not\derives_{\mathsf{IPL}}C\).
In classical ZF define
\[
 S_0=\Gamma,\qquad
 S_{n+1}=
 \begin{cases}
 S_n\cup\{D_n\},&S_n\cup\{D_n\}\not\derives_{\mathsf{IPL}}C,\\
 S_n,&S_n\cup\{D_n\}\derives_{\mathsf{IPL}}C.
 \end{cases}
\]
Then \(T=\bigcup_{n\in\mathbb N}S_n\) is maximal avoiding \(C\),
and hence is the prime extension required above. The supplied enumeration
is the only countability assumption; \(\Gamma\) is still an arbitrary
subset of that formula set.
\end{hblemma}

This enumeration considers every formula for addition whether or not it is
already derivable. A saturation method that scans a fixed disjunction list
and acts only when the current disjunction is derivable must revisit
disjunctions that become derivable later. A modern primary account
of that construction and its countable Hilbert presentation is
\cite[Sections~2--3]{GuoChenBentzen2023IPL}; the LJ proof here gives the
needed derivation operations explicitly.

\begin{hblemma}{thm:p08-ipl-failed-implication}{proved=annex:p8.ipl.complete.failed-implication}{Extension witnessing a failed implication}
If \(T\) is a prime theory and \(A\to B\notin T\), there is a
prime theory \(U\) such that
\[
 T\subseteq U,\qquad A\in U,\qquad B\notin U.
\]
This uses the arbitrary-language prime-extension lemma, or its enumerated
version when an enumeration is supplied.
\end{hblemma}

\subsection{The canonical countermodel}
\label{sec:p08-ipl-canonical-countermodel}

\begin{hbdefinition}{def:p08-ipl-canonical-cone}{Canonical model above a prime theory}
Fix a prime theory \(T_0\). Define
\[
 \begin{split}
 W_{T_0}&=\{T\subseteq\mathcal F:T\text{ is prime and }T_0\subseteq T\},\\
 T\leq U&\quad\Longleftrightarrow\quad T\subseteq U,\\
 V_{T_0}(p)&=\{T\in W_{T_0}:p\in T\}.
 \end{split}
\]
The associated model is \(\mathcal M_{T_0}=(W_{T_0},\subseteq,V_{T_0})\).
It contains every prime extension of its root, including extensions that
contain a formula omitted by \(T_0\).
\end{hbdefinition}

\begin{hblemma}{thm:p08-ipl-canonical-truth}{proved=annex:p8.ipl.complete.truth}{Canonical frame, persistence, and truth lemma}
The canonical model is a set-sized persistent model on a partial order
rooted at \(T_0\). For every \(T\in W_{T_0}\) and every
\(A\in\mathcal F\),
\[
 \mathcal M_{T_0},T\forces A
 \quad\Longleftrightarrow\quad A\in T
 \quad\Longleftrightarrow\quad T\derives_{\mathsf{IPL}}A.
\]
Forcing every formula persists under inclusion of prime theories.
\end{hblemma}

\begin{hbtheorem}{thm:p08-ipl-strong-completeness-full}{proved=annex:p8.ipl.complete.strong}{Soundness and strong completeness of IPL}
In classical ZFC, for every set \(\mathsf{Prop}\), every
\(\Gamma\subseteq\mathcal F\), and every \(A\in\mathcal F\),
\[
 \Gamma\derives_{\mathsf{IPL}}A
 \quad\Longleftrightarrow\quad
 \Gamma\entails_{\mathsf{IPL}}A.
\]
If derivability fails, there is a set-sized partially ordered model with a
root forcing every member of \(\Gamma\) and failing to force \(A\).
The equivalence holds separately for consequence over all partial orders,
over all preorders, and over all rooted partial orders. With a supplied
formula enumeration, these statements have the classical ZF proof
specified above.
\end{hbtheorem}

\begin{hbcorollary}{thm:p08-ipl-semantic-compactness}{proved=annex:p8.ipl.complete.compactness}{Semantic compactness}
Under the hypotheses of Theorem~\ref{thm:p08-ipl-strong-completeness-full},
if \(\Gamma\entails_{\mathsf{IPL}}A\), some finite
\(\Gamma_f\subseteq\Gamma\) already entails \(A\).
Moreover, a premise set is forced at a world of a model exactly when
every finite subset is forced at a world of some model. The witnessing
model for the whole set can be a rooted partial order.
\end{hbcorollary}

\subsection{Finite filtration and explicit countermodel bounds}
\label{sec:p08-ipl-finite-filtration}

\begin{hbdefinition}{def:p08-ipl-filtration}{Filtration by subformula truth sets}
Let \(\mathcal M=(W,\leq,V)\) be a nonempty preorder model and let
\(\Sigma\subseteq\mathcal F\) be a finite set closed under taking
subformulas. In this subsection \(\Sigma\) has precisely this finite
formula-set role. Define
\[
 \tau_\Sigma(w)=\{B\in\Sigma:\mathcal M,w\forces B\},\qquad
 w\sim_\Sigma v\quad\Longleftrightarrow\quad
 \tau_\Sigma(w)=\tau_\Sigma(v).
\]
Write \([w]_\Sigma\) for the equivalence class of \(w\) and put
\[
 \begin{split}
 W_\Sigma&=W/{\sim_\Sigma},\\
 [w]_\Sigma\leq_\Sigma[v]_\Sigma
   &\quad\Longleftrightarrow\quad
     \tau_\Sigma(w)\subseteq\tau_\Sigma(v),\\
 V_\Sigma(p)&=
 \begin{cases}
  \{[w]_\Sigma:w\in V(p)\},&p\in\Sigma,\\
  \varnothing,&p\notin\Sigma.
 \end{cases}
 \end{split}
\]
Let \(\mathcal M_\Sigma=(W_\Sigma,\leq_\Sigma,V_\Sigma)\).
The order compares all the truth information in \(\Sigma\), including
its implication formulas. It can relate classes whose chosen
representatives were incomparable in the original frame.
\end{hbdefinition}

\begin{hblemma}{thm:p08-ipl-filtration-truth}{proved=annex:p8.ipl.complete.filtration}{Filtration lemma, including implication}
The definitions of \(\leq_\Sigma\) and \(V_\Sigma\) are independent
of representatives and make \(\mathcal M_\Sigma\) a persistent model
on a finite partial order. It has at most \(2^{|\Sigma|}\) worlds,
and for all \(B\in\Sigma\) and \(w\in W\),
\[
 \mathcal M_\Sigma,[w]_\Sigma\forces B
 \quad\Longleftrightarrow\quad\mathcal M,w\forces B.
\]
If \(r\) is a root of the original frame, \([r]_\Sigma\) is a root
of the filtered frame.
\end{hblemma}

\begin{hbtheorem}{thm:p08-ipl-finite-rooted-countermodel}{proved=annex:p8.ipl.complete.finite-model}{Finite rooted countermodels}
If \(\Gamma\) is finite and
\(\Gamma\not\derives_{\mathsf{IPL}}A\), let \(\Sigma\) be the
set of all subformulas of formulas in \(\Gamma\cup\{A\}\).
There is a persistent model on a finite partial order with at most
\(2^{|\Sigma|}\) worlds and a root \(r\) such that
\[
 \mathcal M,r\forces\Gamma,
 \qquad \mathcal M,r\not\forces A.
\]
Here negation abbreviations are expanded before subformulas are collected.
The existence proof has the same explicit metatheoretic hypotheses as
Theorem~\ref{thm:p08-ipl-strong-completeness-full}.
\end{hbtheorem}

\begin{hbcorollary}{thm:p08-ipl-finite-input-decision}{proved=annex:p8.ipl.complete.decision}{A terminating decision procedure for finite inputs}
For a finite list of premises \(\Gamma\) and a conclusion \(A\),
finite model search up to the preceding bound decides
\(\Gamma\derives_{\mathsf{IPL}}A\). The input formulas must have a
finite effective encoding with decidable equality of their atom names.
Only the finitely many atoms occurring in the input need valuations.
This is an existence and termination bound for a decision procedure;
it asserts no optimal complexity bound.
\end{hbcorollary}

\begin{hbexample}{def:p08-ipl-completeness-two-worlds}{A root can leave a proposition undecided}
Let the frame be \(r<s\), with reflexive loops, and put
\(V(p)=\{s\}\). At \(r\), the atom \(p\) is not forced, and
\(\neg p\) is not forced because \(s\forces p\). Thus
\(r\not\forces p\lor\neg p\). At \(s\), \(p\) is forced.
Neither world forces \(\neg p\), so both force \(\neg\neg p\).
The root therefore also refutes \(\neg\neg p\to p\): the root itself
is an accessible witness to the failed implication. Soundness establishes
the nonderivability of both formulas. Failure to force a formula and
forcing its negation are different conditions.
\end{hbexample}

\begin{hbexample}{def:p08-ipl-completeness-persistent-premise}{Using an implication at a later world}
Let \(p,q,t\) be atoms; the letter \(t\) has this local role.
The premise \(p\to q\) entails
\((q\to t)\to(p\to t)\). Indeed, let \(w\forces p\to q\).
If \(v\geq w\) forces \(q\to t\), and \(u\geq v\) forces
\(p\), then \(u\geq w\), so \(u\forces q\). Since \(u\geq v\),
the implication forced at \(v\) gives \(u\forces t\). Hence
\(v\forces p\to t\), proving the displayed entailment at \(w\).
The same derivability follows by two applications of the deduction
equivalence and modus ponens.
\end{hbexample}
